\vspace{-0.2cm}

\section{Introduction}

AI for software engineering has made remarkable progress, becoming a notable success within generative AI. Despite this, achieving fully automated software engineering remains a significant challenge, requiring significant research efforts across academia and industry. In this paper, we provide an opinionated view of the tasks, challenges, and promising directions towards achieving this goal.

Many existing surveys focus on particular topics discussed in this work. \citet{liang2024large} and \citet{sergeyuk2025using} survey the successes and challenges of AI programming assistants, \citep{wang2024software} survey using LLMs for software testing, and \citet{joel2024survey} survey using LLMs in low-resource and domain-specific languages, and \citet{zhang2023survey} focus on automated program repair, both with and without LLMs. Finally, \citet{yang2024formal} is a roadmap for formal mathematical reasoning and has some overlap with our discussion on software verification. In addition, many works discuss the current state, challenges, and future of AI for software engineering \citep{fan2023large, ozkaya2023application, wong2023natural, zheng2023survey, hou2024large, jin2024llms, wan2024deep}. Our work draws inspiration from them, and we recommend that the reader consult them for alternative perspectives. 

In this position paper, our goal is threefold. In Sec. \ref{sec:tasks-milestones}, we provide a structured way to think about progress in the field and list tasks important in AI for software engineering. We also provide three measures for categorizing concrete realizations of the task: the scale of the problem, the amount of algorithmic complexity, and the level of human intervention.  In Sec. \ref{sec:challenges}, we highlight nine challenges in the field that today's models face, each of which applies to several tasks. In Sec. \ref{sec:paths}, we posit five research themes that are promising to tackle the aforementioned challenges. We hope that through our paper, the reader can appreciate the progress the field has made, understand the shortcoming of today's state-of-the-art models, and gain concrete research ideas for tackling these challenges.


% - AI coding has made progress last few years + is one of the success stories for genai

% - still a lot to go in terms of getting to fully automated swe

% - goal is to highlight the open problems and provide structure for the field going fwd

% - metrics, tasks, challenges framework

% - 2 purposes for paper existing: 1) highlight everything left to do, 2) provide this framework to structure the field

% alternative views: 
% - AI for code is a solved problem
% - we shouldn't be doing this because of economic reasons
% - AI safety: dangerous to have machines that can program themselves

